\documentclass{UAmathtalk}

\author{Be'eri Greenfeld}
\address{Hunter College, CUNY}
\urladdr{https://sites.google.com/view/beeri-greenfeld}
\title{Noncommutative Points}
\date{Thursday, November 12, 2026}

\renewcommand*{\where}{Massry B008}


\begin{document}

\maketitle

\begin{abstract}
A projective scheme can be recovered as the moduli space of point modules over its homogeneous coordinate ring. This is an important starting point for the rich, yet mysterious, theory of noncommutative projective algebraic geometry.

We connect this theory with symbolic dynamics by studying the moduli spaces associated with monomial algebras that play the role of homogeneous coordinate rings of subshifts. We explore the surprising traces left by dynamics and classical algebraic geometry on these spaces, which can themselves be seen as symbolic pro-algebraic varieties in the sense of Gromov.

In particular, we describe “symbolic projective lines” and exhibit a non-trivial analogy with their quantum counterparts: both families can be described in terms of group algebras of nilpotent groups (of different classes). Finally, we discuss the interaction between nilpotent elements and point modules in light of the projective Nullstellensatz, and present a recent application to the construction of a counterexample to the Köthe Conjecture.
\end{abstract}

\end{document}
